\section{Estimator Model}
SimpleSFT treats a training step as three candidate peaks:
\[
M_{\mathrm{peak}} =
\max(M_{\mathrm{forward}}, M_{\mathrm{backward}}, M_{\mathrm{optimizer}}).
\]
Forward starts from parameters and runtime support. Backward adds gradients and
backward-live activation state. The optimizer step adds update scratch,
communication buckets, and phase carryover. This breakdown gives a useful
structure for model design and calibration because each estimator change can be
attached to a specific phase and memory family.

\subsection{Inputs}
The estimator combines a \texttt{ModelSpec} and an \texttt{EstimatorConfig}.
The model spec supplies architecture and tensor metadata: parameter shapes,
linear-layer roles, attention widths, vocabulary size, layer count, and
tensor-parallel layout annotations. The config supplies the training strategy:
tuning mode, optimizer, dtypes, sequence length, GPU budget, attention backend,
checkpointing, LoRA settings, distributed mode, tensor and sequence parallel
choices, runtime-support assumptions, and DDP/ZeRO bucket sizes.

The inspected model metadata is highly detailed. For a
Qwen-style model, the estimator needs the number of transformer blocks, the
hidden width, the MLP expansion width, the number of attention heads, the number
of key/value heads, the vocabulary projection shape, and the local role of each
large matrix under tensor parallelism. The same total parameter count can lead
to different memory behavior when the attention projection layout, MLP layout,
or vocabulary tensor is sharded differently.

\begin{center}
\small
\begin{tabular}{@{}p{0.28\linewidth}p{0.62\linewidth}@{}}
\toprule
Estimator input & Why it affects peak memory \\
\midrule
Parameter shapes & Drive parameter bytes, gradient bytes, optimizer state, and
tensor-parallel local ownership. \\
Linear-layer list & Identifies LoRA target modules and adapter parameter
counts. \\
Attention layout & Determines query and key/value widths, attention workspace,
and long-context retained state. \\
Distributed strategy & Changes rank ownership of gradients, optimizer state,
parameters, buckets, and gather windows. \\
Checkpointing flag & Changes retained activation state and adds backward
recompute workspace. \\
\bottomrule
\end{tabular}
\normalsize
\end{center}

\subsection{Resident State}
Resident state is the memory floor on one rank:
\[
R_{\mathrm{base}} =
M_{\mathrm{params}} + M_{\mathrm{opt}} + M_{\mathrm{master}}
+ M_{\mathrm{runtime}} + M_{\mathrm{persist}},
\qquad
R_{\mathrm{grad}} = R_{\mathrm{base}} + M_{\mathrm{grads}}.
\]
Parameters are counted from explicit parameter shapes when available, with
tensor-parallel sharding applied to row, column, and vocabulary-parallel
tensors. ZeRO-3 shards trainable parameters across data-parallel ranks. LoRA
keeps the dense base weights resident and makes the adapter tensors trainable,
so gradients and optimizer state scale with adapter size.

Optimizer memory is optimizer-specific. Adam and AdamW keep two state tensors,
RMSProp and Adagrad keep one, SGD keeps none, and Adafactor uses factored state
for matrix-like tensors. ZeRO-2 and ZeRO-3 shard optimizer state and gradients
for tested optimizer paths. Runtime support is represented with named CUDA,
allocator, NCCL, DeepSpeed, and backend-buffer terms.

The local tensor count is computed before bytes are assigned. Tensor-parallel
metadata determines whether a parameter is replicated or divided across the
tensor-parallel group. ZeRO then applies data-parallel sharding to the training
state owned by its stage. This order preserves the difference between a matrix
partitioned by tensor parallelism and a state object partitioned by ZeRO.

\begin{center}
\small
\begin{tabular}{@{}lll@{}}
\toprule
Object & Full fine-tuning & LoRA fine-tuning \\
\midrule
Base weights & resident & resident \\
Trainable weights & dense model tensors & adapter tensors \\
Gradients & dense trainable tensors & adapter gradients \\
Optimizer state & dense trainable tensors & adapter optimizer state \\
\bottomrule
\end{tabular}
\normalsize
\end{center}

This distinction is central for LLMs, because a full fine-tune must carry
optimizer and gradient state for the dense model. A LoRA run still stores the
base weights during forward and backward passes, while trainable-state memory demand is determined by
the targeted projection layers and the selected adapter rank.

For one dense matrix $W \in \mathbb{R}^{d_{\mathrm{out}}\times d_{\mathrm{in}}}$,
full fine-tuning makes all $d_{\mathrm{out}}d_{\mathrm{in}}$ elements trainable.
LoRA inserts
\[
\Delta W = \frac{\alpha}{r}BA,\qquad
A \in \mathbb{R}^{r\times d_{\mathrm{in}}},\quad
B \in \mathbb{R}^{d_{\mathrm{out}}\times r},
\]
so the trainable count for that target becomes
$r(d_{\mathrm{in}} + d_{\mathrm{out}})$. The dense base matrix remains part of
resident parameter memory, while gradient and optimizer memory follow the
adapter tensors. This gives the estimator a direct way to reason about rank,
target module selection, and optimizer state without introducing a separate
LoRA-only formula for the rest of the model.

Optimizer state is handled at parameter-tensor granularity. Adam-style
optimizers attach two state tensors to each trainable tensor. Adafactor uses
row and column factors for matrix-like tensors, which is materially smaller for
large projection matrices. This tensor-level treatment matters for mixed
training modes: a run can have dense frozen base parameters, trainable adapter
matrices, and optimizer state only for those adapters.

\subsection{Retained Activations and Attention}
Retained activation memory is a sum of terms:
\[
M_{\mathrm{act}} =
M_{\mathrm{visible}} + M_{\mathrm{saved\ inputs}} + M_{\mathrm{mlp}}
+ M_{\mathrm{attn}} + M_{\mathrm{lora}} + M_{\mathrm{ckpt}}
+ M_{\mathrm{loss}} + M_{\mathrm{widened}}.
\]
Full fine-tuning grows with layer count, local token count, hidden width, MLP
width, and attention-path width. LoRA reduces trainable activation pressure by
moving gradients onto targeted adapter paths and low-rank intermediates.
Gradient checkpointing stores fewer layer internals and shifts work into
backward recomputation.

Attention backend matters most at long context length. Standard attention
requires attention-score states that scale quadratically with sequence length.
FlashAttention removes the need to store quadratic-length attention scores by
using tiled computation and recomputation, with a smaller additional workspace
term \cite{dao2022flashattention}. SimpleSFT therefore models FlashAttention
primarily through retained attention-path context and workspace terms.

Gradient checkpointing is modeled as a change to the activation life cycle.
Layer internals that would be retained until backward are reduced, while a
separate recompute window is added during backward. The estimator represents
this through checkpoint-boundary bytes, saved linear-input bytes, residual and
normalization context, and recompute workspace. This keeps the memory tradeoff
visible: checkpointing lowers retained forward state and raises backward
compute plus transient workspace.

\subsection{Workspace and Phase Assembly}
Workspace is modeled as a live window. The main terms are attention workspace,
backward kernel workspace, recompute workspace, loss workspace, tensor/sequence
parallel communication windows, optimizer update scratch, and distributed
bucket objects.

DDP contributes a reducer bucket and a communication-overlap window. ZeRO
contributes all-gather, reduce, and prefetch buckets. Bucket capacities are
derived from configured element counts and clipped by the local payload that
can fill them. ZeRO optimizer pressure is the maximum of fetch, update, and
communication subphases.

The bucket model uses configured capacities and local payload caps. A reducer
or all-gather bucket can reserve a large object, yet the active payload remains
bounded by the data held on that rank. This makes the estimate responsive to
both runtime settings and model size. ZeRO-2 keeps parameters resident while
sharding optimizer state and gradients. ZeRO-3 also shards trainable
parameters, adding parameter gather windows around computation.

The final phase equations are:
\[
\begin{aligned}
M_{\mathrm{forward}} &= R_{\mathrm{base}} + A_{\mathrm{fwd}} + W_{\mathrm{fwd}}, \\
M_{\mathrm{backward}} &= R_{\mathrm{grad}} + A_{\mathrm{bwd}} + W_{\mathrm{bwd}}, \\
M_{\mathrm{optimizer}} &=
R_{\mathrm{grad}} + E_{\mathrm{bwd}} + W_{\mathrm{opt}} + C_{\mathrm{opt}}.
\end{aligned}
\]
Here $E_{\mathrm{bwd}}$ is backward-end state such as gradients and live
distributed buckets, and $C_{\mathrm{opt}}$ represents carryover that remains
live at optimizer-step entry.

The phase assembly makes the peak source explicit. A forward peak usually
signals parameter plus retained activation pressure. A backward peak usually
signals the combination of gradients, retained or recomputed activations, and
communication workspace. An optimizer-step peak usually signals optimizer
state, update scratch, and ZeRO communication windows. Reporting this phase is
important because two configurations can have similar total peaks while needing
different fixes. For example, checkpointing helps a forward-activation peak,
while ZeRO stage changes help persistent optimizer or gradient pressure.

\begin{center}
\small
\begin{tabular}{@{}p{0.22\linewidth}p{0.68\linewidth}@{}}
\toprule
Peak phase & Dominant terms SimpleSFT reports \\
\midrule
Forward & Base resident state, retained activations, attention workspace,
loss-related state. \\
Backward & Resident state with gradients, backward activation state,
checkpoint recompute workspace, DDP or ZeRO communication windows. \\
Optimizer step & Gradients, optimizer state, update scratch, ZeRO fetch and
communication windows, carried state from backward. \\
\bottomrule
\end{tabular}
\normalsize
\end{center}
